\lecture{39}{Крупномасштабное машинное обучение}

\sourcevideo
    {https://www.youtube.com/playlist?list=PLOzRYVm0a65fhKDS8cYIC7YCuVGJ4FOJx}
    {Week 10: Lecture 39: Large Scale Machine Learning}

Пусть параметр модели равен \(w\in\mathbb R^d\), а данные распределены
между \(N\) вычислительными работниками. При большой размерности
модели передача параметров и градиентов может занимать больше времени,
чем локальное вычисление. Ниже рассматриваются две архитектуры
синхронного обучения --- сервер параметров и кольцевой all-reduce, ---
а затем приводится базовая оценка для стохастического градиентного
спуска.

\section{Распределённый градиентный шаг}

Работник \(i\) хранит локальный набор \(\mathcal D_i\) мощности
\(n_i=\lvert\mathcal D_i\rvert\). При
\[
    M=\sum_{i=1}^{N}n_i
\]
глобальная и локальные эмпирические функции имеют вид
\[
    F(w)
    =
    \frac1M
    \sum_{i=1}^{N}
    \sum_{\zeta\in\mathcal D_i}
    \ell(w;\zeta),
    \qquad
    F_i(w)
    =
    \frac1{n_i}
    \sum_{\zeta\in\mathcal D_i}
    \ell(w;\zeta).
\]
Следовательно,
\[
    F(w)
    =
    \sum_{i=1}^{N}\frac{n_i}{M}F_i(w),
    \qquad
    \nabla F(w)
    =
    \sum_{i=1}^{N}\frac{n_i}{M}\nabla F_i(w).
\]
Данные остаются у работников; для общего шага требуется агрегировать
только локальные градиенты.

На итерации \(k\) все работники используют один параметр \(w^k\) и
вычисляют
\[
    g_i^k=\nabla F_i(w^k).
\]
После агрегации
\[
    g^k
    =
    \sum_{i=1}^{N}\frac{n_i}{M}g_i^k,
    \qquad
    w^{k+1}=w^k-\eta_k g^k.
\]
При точных локальных градиентах имеем \(g^k=\nabla F(w^k)\), поэтому
это обычный централизованный градиентный шаг, распределённый между
несколькими устройствами.

\section{Сервер параметров}

В архитектуре сервера параметров центральный узел хранит \(w^k\),
рассылает его работникам, принимает локальные градиенты, вычисляет их
взвешенную сумму и формирует \(w^{k+1}\). Получаемые векторы можно
сразу добавлять к накопленной сумме, поэтому без учёта хранения самой
модели дополнительная память сервера имеет порядок \(O(d)\).

За один раунд сервер принимает \(N\) векторов градиента и отправляет
\(N\) копий параметра. Его коммуникационная нагрузка имеет порядок
\[
    O(Nd).
\]
При росте \(N\) центральный узел становится узким местом и единой
точкой отказа. Преимущество этой архитектуры состоит в простоте и
малом числе логических стадий одного глобального шага.

\section{Кольцевой all-reduce}

Для синхронного шага каждому работнику достаточно получить сумму
локальных вкладов. При равных размерах выборок это
\[
    G=\sum_{i=1}^{N}g_i,
    \qquad
    w^{k+1}=w^k-\eta_k\frac1N G.
\]
При различных \(n_i\) работник заранее умножает свой градиент на
\(n_i/M\), после чего all-reduce суммирует уже взвешенные векторы.

В кольцевой реализации каждый участник передаёт сообщения следующему
работнику и получает их от предыдущего. Предположим сначала, что
\(d\) делится на \(N\), и разобьём
\[
    g_i
    =
    \bigl(g_i^{(0)},\ldots,g_i^{(N-1)}\bigr),
    \qquad
    g_i^{(r)}\in\mathbb R^{d/N}.
\]
При неделимости используются блоки близких размеров или дополнение
нулями. На блоки делится передаваемый вектор, а не локальный набор
данных.

Первая фаза, scatter-reduce, длится \(N-1\) раундов. В каждом раунде
участник отправляет один блок вперёд по кольцу, принимает другой блок
и прибавляет его к соответствующей локальной сумме. После завершения
фазы каждый работник хранит один полностью просуммированный блок
\[
    G^{(r)}=\sum_{i=1}^{N}g_i^{(r)},
\]
а разные блоки распределены между разными работниками.

Например, для
\[
    g_0=(A_0,B_0,C_0),\qquad
    g_1=(A_1,B_1,C_1),\qquad
    g_2=(A_2,B_2,C_2)
\]
после двух раундов получены суммы
\[
    A=A_0+A_1+A_2,\qquad
    B=B_0+B_1+B_2,\qquad
    C=C_0+C_1+C_2,
\]
причём каждая из них пока находится у одного участника.

Во второй фазе, all-gather, готовые блоки циркулируют по кольцу без
дальнейшего суммирования. Ещё через \(N-1\) раундов каждый работник
получает полный вектор
\[
    G
    =
    \bigl(G^{(0)},\ldots,G^{(N-1)}\bigr)
    =
    \sum_{i=1}^{N}g_i
\]
и независимо выполняет то же обновление параметра.

\section{Коммуникационная стоимость}

В одном раунде передаётся примерно \(d/N\) компонент. Поэтому за
scatter-reduce один работник передаёт \((N-1)d/N\) компонент и столько
же --- за all-gather. Полный объём передачи на одного работника равен
\[
    2\frac{N-1}{N}d
    \longrightarrow
    2d
    \qquad
    \text{при }N\to\infty.
\]
Таким образом, объём данных на одного участника имеет порядок
\(O(d)\), но один глобальный шаг требует \(2(N-1)\) последовательных
коммуникационных раундов. Ring all-reduce устраняет центральный
агрегатор и равномерно распределяет трафик, однако его задержка может
расти с числом работников.

Сервер параметров выгоден малым числом стадий, но концентрирует
нагрузку порядка \(Nd\) в одном узле. Кольцо распределяет нагрузку и
оставляет на каждом работнике объём порядка \(d\), но требует
линейного числа раундов. Выбор зависит от размерности модели, числа
работников, пропускной способности, задержки и надёжности сети.
Путь также имеет малую степень вершин, но большой диаметр; звезда
имеет диаметр два, но перегружает центр. Статические экспоненциальные
графы позволяют получить промежуточный компромисс между степенью и
диаметром.

Обе рассмотренные архитектуры предполагают, что до коллективной
операции все работники имеют одинаковый \(w^k\). После получения
одного агрегированного градиента они снова приходят к общему
\(w^{k+1}\). Поэтому эти схемы распределяют вычисление
централизованного шага, но не согласуют изначально различные локальные
модели.

\section{Стохастический градиентный спуск}

Пусть на итерации \(k\) используется случайный градиент
\[
    g_k=g(w^k;\zeta_k),
    \qquad
    \mathbb E[g_k\mid w^k]=\nabla F(w^k),
\]
и его условная дисперсия ограничена:
\[
    \mathbb E\!\left[
        \lVert g_k-\nabla F(w^k)\rVert^2
        \,\middle|\,
        w^k
    \right]
    \le
    \sigma^2.
\]
Для обновления \(w^{k+1}=w^k-\eta g_k\), функции с
\(L\)-липшицевым градиентом и нижней гранью \(F_{\inf}\), при
\(0<\eta\le 1/L\) выполняется стандартная оценка
\[
    \frac1T
    \sum_{k=0}^{T-1}
    \mathbb E\lVert\nabla F(w^k)\rVert^2
    \le
    \frac{2\bigl(F(w^0)-F_{\inf}\bigr)}{\eta T}
    +
    L\eta\sigma^2.
\]
Первое слагаемое уменьшается при росте шага, а второе отражает влияние
шума. При \(\eta\asymp T^{-1/2}\) получаем
\[
    \frac1T
    \sum_{k=0}^{T-1}
    \mathbb E\lVert\nabla F(w^k)\rVert^2
    =
    O\!\left(T^{-1/2}\right).
\]
Константа зависит от начального функционального разрыва, гладкости и
дисперсии. Для невыпуклой функции эта оценка характеризует
приближение к стационарности, а не к глобальному минимуму. Чтобы
средний квадрат нормы градиента не превосходил \(\varepsilon\), базовая
оценка требует \(T=O(\varepsilon^{-2})\) итераций.

\section{Переход к децентрализованному SGD}

Сервер параметров и ring all-reduce воспроизводят один
централизованный SGD-шаг. В полностью децентрализованной схеме агент
\(i\) хранит собственную переменную \(w_i^k\) и обменивается только с
соседями. Тогда к стохастической ошибке градиента добавляются
рассогласование локальных моделей и зависимость скорости смешивания
от спектра коммуникационного графа. Эти эффекты рассматриваются в
следующей лекции.
